% ----------------------
\subsection{Section 4 (February 1829)}
% ----------------------
	\label{DC4}
	
	% -----
	\subsubsection{Historical Background}
	% -----
		\label{DC4-HB}
		
		% --
		\paragraph{JSP}
		% --
		
			``JS dictated this revelation for his father, Joseph Smith Sr., one of his earliest and staunchest supporters. When John Whitmer copied this text into Revelation Book 1, he included this heading: `A Revelation to Joseph the Father of the Seer he desired to know what the Lord had for him to do \& this is what he Received as follows.' Revelation Book 1 initially gave the date of 1828. An unidentified scribe wrote a `9' over the `8,' thus changing the date from 1828 to 1829, apparently correcting a scribal error. The index to Revelation Book 1 also lists 1829 as the date of the revelation. Sidney Rigdon, likely in late 1831, added `Fe\uline{b}\supersu{r.}' to the heading in Revelation Book 1 to further specify the date. The copy featured below is a more complete and probably an earlier version than that inscribed in Revelation Book 1, which is missing the page that includes the final portion of this revelation. The version below is in the handwriting of Edward Partridge and was kept by him. Partridge dated the document to 1829, a date also used in JS's history.
			
			``Joseph Knight Sr., another early supporter of JS, wrote that Joseph Sr. and Samuel Smith stopped at his home in Colesville, New York, in January 1829 before going on to visit JS and Emma Smith. `I told him [Joseph Smith Sr.] they had traviled far enough,' Knight wrote, `[and] I would go with my sley and take them Down [to Harmony] to morrow[.] I went Down and found them well and the[y] were glad to see us[.] we conversed about many things. in the morning I gave the old man a half a Dollar and Joseph a little money to Buoy [buy] paper to translate.' JS had apparently not translated since June 1828, and Knight's provision of paper may have allowed him to resume translation. Within weeks of Knight's visit, JS began translating again, with Emma, Samuel, and Martin Harris each acting briefly as scribe.
			
			``Dictated shortly before the translation work resumed, this revelation spoke of a `marvelous work' about to come forth and added that the `field is white already to harvest.' These phrases, also used in several JS revelations in the spring of 1829, invoked a sense of urgency and an impending spiritual harvest. Though addressed to Joseph Smith Sr., this revelation was written as if it could apply to all who read it.
			
			``The degree to which Joseph Smith Sr. acted upon this revelation is unknown, but his call `to the work' may have had a significant immediate impact when he returned to Palmyra, New York, where Oliver Cowdery was boarding at his house. Joseph Sr. and Lucy Mack Smith had met Cowdery when he began teaching school in the Manchester, New York, district late in the fall of 1828. Lucy wrote that although Cowdery had questioned Joseph Sr. about the gold plates, he `did not succeed in eliciting any information' for `a long time.' This revelation may have prompted Joseph Sr. to share a `sketch of the facts which related to the plates' with Cowdery, who became convinced that he had been called by God to assist JS as his scribe'' (D1:10--13).
		
	% -----
	\subsubsection{Versions}
	% -----
		\label{DC4-V}
	
		\begin{tabular}{p{1in}p{2in}p{1in}p{2in}}
			JSP		& \hyperref[EMS-1828-1829-JS-FEB-1829]{February 1829}	& BC		& Chapter 3, 9 \\
			DC 1835 & Section 31, 158										& TS		& 3:16 (15 June 1842), 817 \\
			MS		& 3:7 (November 1842), 118								& DC 1844	& Section 31, 231--232 \\
		\end{tabular}
		
		\medskip
		
		See the \href{https://drive.google.com/file/d/1goAvNz8jS4R8slYfMG_db55Ty2z_vmgK/view?usp=sharing}{sections comparison} document for more information about the version history.

	% -----	
	\subsubsection{Section 4:1} % *DC4-1 
	% -----
		\label{DC4-1}

			\footnote{%
				% \label{}%
				The JSP version, which is from a loose leaf, has this heaing: ``A revalation from the Lord unto Jos [\textit{page damaged}] AD 1829 Saying.'' RB1 has, ``AD 1828 A Revelation to Joseph [Smith Sr.] the Father of the Seer he desired to know what the Lord had for him to do \& this is what he Received as follows Saying.'' BC has, ``1 \textit{A Revelation given to Joseph [Smith Sr.], the father of Joseph, in Harmony, Pennsylvania, February, 1829. saying:}''. DC 1835 has, ``\textit{Revelation to Joseph Smith, Sen., given February, 1829.}'' DC 1844 has, ``\textit{Revelation to Joseph Smith, Sen. given February, 1829.}''
				}%
		NOW behold, a marvelous work is about to come forth among the
			\footnote{%
				% \label{}%
				``Children of men'' is an OT term; it doesn't appear in the KJV OT. See notes at \hyperref[GENESIS-11-5]{Genesis 11:5}. The term seems intended to invoke the difference between us and God.
				}%
		children of men.

		\begin{framed}
		\scriptsize
			\begin{compactdesc}
				\item[JSP] now behold a marvelous work is about to come among the [\textit{page damaged}] [children] of men
				\item[RB1] now Behold a Marvelous work is about to come forth among the children of men 
				\item[BC] NOW, behold, a marvelous work is about to come forth among the children of men, 
				\item[DC 1835] 1 Now, behold, a marvellous work is about to come forth among the children of men,
				\item[DC 1844] 1 Now, behold, a marvellous work is about to come forth among the children of men,
			\end{compactdesc}
		\end{framed}

	% -----
	\subsubsection{Section 4:2} % *DC4-2 
	% -----
		\label{DC4-2}

		Therefore, O ye that embark%
			\footnote{%
				% \label{}%
				This is the only occurrence of ``embark'' in all the scriptures.
				}
		in the service of God, see that ye serve him with all your heart, might, mind and strength,%
			\footnote{%
				% \label{}%
				This combination of four things only appears two other times in the scriptures---at \hyperref[DC59-5]{DC 59:5} and \hyperref[DC98-47]{DC 98:47}. There is a bit of biblical precursor at \hyperref[MARK-12-30]{Mark 12:30} and \hyperref[LUKE-10-27]{Luke 10:27}. Notice that the biblical precursors are both passages about loving God, though the one in Luke is spoken by a lawyer, and not by Christ. See the OT references at the Luke passage. See also \hyperref[DC64-22]{DC 64:22} and \hyperref[DC64-34]{34}; \hyperref[MORONI-7-13]{Moroni 7:13};
				}
		that ye may stand blameless before God at the last day.

		\begin{framed}
		\scriptsize
			\begin{compactdesc}
				\item[JSP] Therefore O ye that embark in the service of God see that ye serve him with all your heart \sout{mind} might mind \& strenght that ye may stand blameless before God at the last day
				\item[RB1] therefore O ye that embark in the service of God see that ye Serve him with all your heart might mind \& Strength that ye may stand blameless before God at the last day 
				\item[BC] therefore, O ye that embark in the service of God, see that ye serve him with all your heart, might, mind and strength, that ye may stand blameless before God at the last day:
				\item[DC 1835] therefore, O ye that embark in the service of God, see that ye serve him with all your heart, might, mind and strength, that ye may stand blameless before God at the last day:
				\item[DC 1844] therefore, O ye that embark in the service of God, see that ye serve him with all your heart, might, mind and strength, that ye may stand blameless before God at the last day:
			\end{compactdesc}
		\end{framed}

	% -----
	\subsubsection{Section 4:3} % *DC4-3 
	% -----
		\label{DC4-3}

		Therefore, if ye have desires to serve God ye are called to the work;%
			\footnote{%
				% \label{}%
				The way this is given here, it's given as a general principle, not restricted to people who are involved in a restoration faith. I think you see this born out among other people who have desires to serve God and so feel ``called'' to the work. Compare \hyperref[DC6-4]{DC 6:4}, which says that those who are already working are ``called of God.'' 
				}

		\begin{framed}
		\scriptsize
			\begin{compactdesc}
				\item[JSP] therefore if ye have desires to serve God ye are called to the work
				\item[RB1] therefore if ye have desires to serve God ye are Called to the [\textit{pages 3--10 missing}]
				\item[BC] Therefore, if ye have desires to serve God, ye are called to the work, 
				\item[DC 1835] therefore, if ye have desires to serve God, ye are called to the work,
				\item[DC 1844] therefore, if ye have desires to serve God, ye are called to the work,
			\end{compactdesc}
		\end{framed}

	% -----
	\subsubsection{Section 4:4} % *DC4-4 
	% -----
		\label{DC4-4}

		For behold the field is white already to harvest; and lo, he that thrusteth in his sickle with his might, the same layeth up in store that he perisheth not, but bringeth salvation to his soul;%
			\footnote{%
				% \label{}%
				See \hyperref[DC61-33]{DC 61:33--34}.
				}

		\begin{framed}
		\scriptsize
			\begin{compactdesc}
				\item[JSP] for behold the field is white already to harvest \& lo he that thursteth in his sickle with his might the same layeth up his store that he perish not but bringeth Salvation to his \sout{own} soul
				\item[RB1] [\textit{missing}]
				\item[BC] for behold, the field is white already to harvest, and lo, he that thrusteth in his sickle with his might, the same layeth up in store that he perish not, but bringeth salvation to his soul, 
				\item[DC 1835] for behold the field is white already to harvest, and lo, he that thrusteth in his sickle with his might, the same layeth up in store that he perish not, but bringeth salvation to his soul,
				\item[DC 1844] for behold the field is white already to harvest, and lo, he that thrusteth in his sickle with his might, the same layeth up in store that he perish not, but bringeth salvation to his soul,
			\end{compactdesc}
		\end{framed}

	% -----
	\subsubsection{Section 4:5} % *DC4-5 
	% -----
		\label{DC4-5}

		And faith, hope, charity and love,%
			\footnote{%
				% \label{}%
				Note that ``charity'' and ``love'' are mentioned separately here. This is pretty unusual---``charity'' is just identified with ``love'' elsewhere, including \hyperref[2NEPHI-26-30]{2 Nephi 26:30}, \hyperref[ETHER-12-34]{Ether 12:34}, \hyperref[MORONI-7-47]{Moroni 7:47}, and \hyperref[MORONI-8-17]{Moroni 8:17}. You do get the pair occurring together (and not identified with each other) again at \hyperref[DC12-8]{DC 12:8}, which is dated to May 1829. You might think of ``love'' here as \gr{φιλία}, though that word only occurs a single time in the NT, at \hyperref[JAMES-4-4]{James 4:4}. Of course that word is very important outside the NT.
				}
		with an eye \hyperref[SINGLE]{single} to the glory of God, qualify%
			\footnote{%
				% \label{}%
				Here is ``qualify'' in the 1828 dictionary:
					\begin{compactitem}
						\item[1] To fit for any place, office, occupation or character; to furnish with the knowledge, skill or other accomplishment necessary for a purpose.
						\item[2] To make capable of any employment or privilege; to furnish with legal power or capacity.
						\item[3] To abate; to soften; to diminish.
						\item[4] To ease; to assuage.
						\item[5] To modify; to restrain; to limit by exceptions.
						\item[6] To modify; to regulate; to vary.
					\end{compactitem}
				So the first two definitions are the ones you want, I think.
				}
		him for the work.

		\begin{framed}
		\scriptsize
			\begin{compactdesc}
				\item[JSP] \& faith hope charity \& love with an eye single to the glory of God constitutes him for the work
				\item[RB1] [\textit{missing}]
				\item[BC] and faith, hope, charity, and love, with an eye single to the glory of God, qualifies him for the work.
				\item[DC 1835] and faith, hope, charity, and love, with an eye single to the glory of God, qualifies him for the work.
				\item[DC 1844] and faith, hope, charity, and love, with an eye single to the glory of God, qualifies him for the work.
			\end{compactdesc}
		\end{framed}

	% -----
	\subsubsection{Section 4:6} % *DC4-6 
	% -----
		\label{DC4-6}

		Remember faith, virtue, knowledge, temperance, patience, brotherly kindness, godliness, charity, humility, diligence.%
			\footnote{%
				% \label{}%
				Compare \hyperref[2PETER-1-4]{2 Peter 1:4--8}.
				}

		\begin{framed}
		\scriptsize
			\begin{compactdesc}
				\item[JSP] remember temperance patience humility diligence \&C.
				\item[RB1] [\textit{missing}]
				\item[BC] 2 Remember temperance, patience, humility, diligence, \&c.,
				\item[DC 1835] 2 Remember faith, virtue, knowledge, temperance, patience, brotherly kindness, godliness, charity, humility, diligence.---
				\item[DC 1844] 2 Remember faith, virtue, knowledge, temperance, patience, brotherly kindness, godliness, charity, humility, diligence.
			\end{compactdesc}
		\end{framed}

	% -----
	\subsubsection{Section 4:7} % *DC4-7 
	% -----
		\label{DC4-7}

			\footnote{%
				% \label{}%
				This seems to be a reference to Christ's instruction at \hyperref[MATTHEW-7-7]{Matthew 7:7} and \hyperref[LUKE-11-9]{Luke 11:9}. This verse doesn't have the part about seeking and finding, so it may actually be drawing from \hyperref[3NEPHI-27-29]{3 Nephi 27:29}, where Christ omits the seeking/finding part and says, ``Therefore ask and ye shall receive; knock and it shall be opened unto you. For he that asketh receiveth; and unto him that knocketh it shall be opened.''
				}%
		Ask, and ye shall receive; knock, and it shall be opened unto you.  Amen.

		\begin{framed}
		\scriptsize
			\begin{compactdesc}
				\item[JSP] Ask \& ye shall receive knock \& it shall be opened unto you amen
				\item[RB1] [\textit{missing}]
				\item[BC] ask and ye shall receive, knock and it shall be opened unto you: Amen.
				\item[DC 1835] Ask and ye shall receive, knock and it shall be opened unto you. Amen.
				\item[DC 1844] Ask and ye shall receive, knock and it shall be opened unto you: Amen.
			\end{compactdesc}
		\end{framed}